\documentclass[10pt, aspectratio=169]{beamer}
% Preámbulo modular para presentaciones Beamer (Modelación Estadística)
% Diseño y paleta institucional del Tecnológico de Monterrey

\documentclass[aspectratio=169,xcolor={dvipsnames,table}]{beamer}

\usepackage[utf8]{inputenc}
\usepackage[spanish,mexico]{babel}
\usepackage{amsmath,amssymb,amsfonts,mathtools}
\usepackage{graphicx}
\usepackage{listings}
\usepackage{booktabs}
\usepackage{tikz}
\usetikzlibrary{matrix,arrows,decorations.pathmorphing,shapes.geometric,positioning}

% Paleta Institucional Tecnológico de Monterrey
\definecolor{TecAzul}{HTML}{0039A6}       % Azul TEC: color primario institucional
\definecolor{TecAzulOscuro}{HTML}{1A2E51} % Azul marino oscuro
\definecolor{TecRojo}{HTML}{EC2661}       % Rojo de acento (paleta miTec)
\definecolor{TecGris}{HTML}{646464}       % Gris neutro para comentarios y subtítulos
\definecolor{TecFondoBloque}{HTML}{F4F6F9}% Fondo suave para bloques

% Configuración de Tema Beamer
\usetheme{Boadilla}
\usecolortheme[named=TecAzulOscuro]{structure}
\setbeamercolor{alerted text}{fg=TecRojo}
\setbeamercolor{palette primary}{bg=TecAzulOscuro,fg=white}
\setbeamercolor{palette secondary}{bg=TecAzul,fg=white}
\setbeamercolor{palette tertiary}{bg=TecAzulOscuro!80!black,fg=white}
\setbeamercolor{title}{fg=TecAzulOscuro,bg=TecFondoBloque}
\setbeamercolor{frametitle}{fg=TecAzulOscuro,bg=TecFondoBloque}

% Bloques personalizados elegantes
\setbeamercolor{block title}{bg=TecAzulOscuro,fg=white}
\setbeamercolor{block body}{bg=TecFondoBloque,fg=black}
\setbeamercolor{block title alerted}{bg=TecRojo,fg=white}
\setbeamercolor{block body alerted}{bg=TecRojo!8,fg=black}
\setbeamercolor{block title example}{bg=TecAzul,fg=white}
\setbeamercolor{block body example}{bg=TecAzul!8,fg=black}

% Redondear bordes y añadir sombra sutil
\setbeamertemplate{blocks}[rounded][shadow=true]
\setbeamertemplate{navigation symbols}{}

% Pie de página limpio con número de diapositiva
\setbeamertemplate{footline}{
  \leavevmode%
  \hbox{%
  \begin{beamercolorbox}[wd=.7\paperwidth,ht=2.25ex,dp=1ex,left]{author in head/foot}%
    \hspace*{2ex}\insertshortauthor~~(\insertshortinstitute) \hfill \insertshorttitle
  \end{beamercolorbox}%
  \begin{beamercolorbox}[wd=.3\paperwidth,ht=2.25ex,dp=1ex,right]{date in head/foot}%
    \insertshortdate{}\hspace*{2em}
    \insertframenumber{} / \inserttotalframenumber\hspace*{2ex}
  \end{beamercolorbox}}%
  \vskip0pt%
}

% Configuración de Listings de Python para visualización pedagógica de código
\lstset{
    language=Python,
    basicstyle=\ttfamily\footnotesize,
    keywordstyle=\color{TecAzulOscuro}\bfseries,
    stringstyle=\color{TecRojo},
    commentstyle=\color{TecGris}\itshape,
    showstringspaces=false,
    breaklines=true,
    frame=lines,
    rulecolor=\color{TecAzulOscuro!30},
    backgroundcolor=\color{TecFondoBloque},
    aboveskip=0.5em,
    belowskip=0.5em,
    literate={á}{{\'a}}1 {é}{{\'e}}1 {í}{{\'i}}1 {ó}{{\'o}}1 {ú}{{\'u}}1
             {Á}{{\'A}}1 {É}{{\'E}}1 {Í}{{\'I}}1 {Ó}{{\'O}}1 {Ú}{{\'U}}1
             {ñ}{{\~n}}1 {Ñ}{{\~N}}1 {¿}{{?`}}1 {¡}{{!`}}1
}

% Metadatos institucionales por defecto
\title[Modelación Estadística]{Modelación Estadística para Ciencia de Datos e Ingeniería}
\author[J. Castillo Colmenares]{Juliho David Castillo Colmenares}
\institute[Tec de Monterrey]{Tecnológico de Monterrey \\ Escuela de Ingeniería y Ciencias}
\date{\today}

% Comandos y macros estadísticas compartidas para las presentaciones Beamer (Metropolis)

% Conjuntos numéricos básicos
\newcommand{\R}{\mathbb{R}}
\newcommand{\N}{\mathbb{N}}
\newcommand{\Q}{\mathbb{Q}}
\newcommand{\Z}{\mathbb{Z}}
\newcommand{\set}[1]{\left\{ #1 \right\}}
\newcommand{\abs}[1]{\left|#1\right|}
\newcommand{\norm}[1]{\left\| #1 \right\|}

% Comandos estadísticos y probabilísticos del curso
\newcommand{\Var}{\operatorname{Var}}
\newcommand{\cov}{\operatorname{Cov}}
\newcommand{\corr}{\rho}
\newcommand{\comb}[2]{\begin{pmatrix} #1 \\ #2 \end{pmatrix}}
\newcommand{\s}{\sigma}
\newcommand{\particion}{\mathfrak{P}}
\newcommand{\card}[1]{\#\left( #1 \right)}
\newcommand{\seq}[3]{\set{#1}_{#2}^{#3}}
\newcommand{\E}{\mathbb{E}}
\newcommand{\Prob}{\mathbb{P}}

% Entornos pedagógicos y cajas en español académico natural
\newenvironment{discusion}{%
  \begin{alertblock}{Pregunta de Discusión}%
}{%
  \end{alertblock}%
}

\newenvironment{reflexion}{%
  \begin{alertblock}{Reflexión para el Aula}%
}{%
  \end{alertblock}%
}

\newenvironment{paradoja}{%
  \begin{alertblock}{Paradoja de Exploración}%
}{%
  \end{alertblock}%
}

\newenvironment{motivacion}{%
  \begin{exampleblock}{Motivación: Ciencia de Datos en la Práctica}%
}{%
  \end{exampleblock}%
}

\newenvironment{retoclase}{%
  \begin{block}{Reto Rápido en Equipo}%
}{%
  \end{block}%
}


\title[Matemáticas de la Regresión]{Sección 09.03: Matemáticas de la Regresión Lineal}
\subtitle{Derivación de MCO, Descomposición SCT=SCR+SCE, y Propiedades de $R^2$}
\author[J. Castillo Colmenares]{Juliho Castillo Colmenares}
\institute{Tecnológico de Monterrey}
\date{\vspace{-1.2cm}}

\begin{document}

% Slide 1: Portada
\begin{frame}[plain]
  \vspace{-0.8cm}
  \titlepage
\end{frame}

% Slide 2: Hoja de Ruta
\begin{frame}{Hoja de Ruta --- Capítulo 09: Regresiones Lineales y Múltiples}
  \scriptsize
  ¿Dónde nos encontramos en el desarrollo modular del capítulo?
  \begin{itemize}\itemsep=0.02em
    \item \textbf{09.01} Correlación como Premisa de la Regresión
    \item \textbf{09.02} Introducción a la Regresión Lineal
    \item \textbf{09.03} Matemáticas de la Regresión: Derivación MCO y $R^2$ \emph{(Hoy)}
    \item \textbf{09.04} Regresión Lineal sobre Datos Simulados
    \item \textbf{09.05} Coeficientes Óptimos, Pruebas $t$/$F$ y RSE
    \item \textbf{09.06} Implementación en Python con \texttt{statsmodels}
    \item \textbf{09.07} Regresión Múltiple, Ecuación Normal y Ridge/Lasso
    \item \textbf{09.08} Validación de Modelos y $k$-fold Cross-Validation
    \item \textbf{09.09} Diagnóstico de Residuos y Multicolinealidad (VIF)
    \item \textbf{09.10} Regresión con \texttt{scikit-learn}
    \item \textbf{09.11} Variables Categóricas y Variables Muda
    \item \textbf{09.12} Transformaciones No Lineales y Regresión Polinomial
  \end{itemize}
  \pause
  \vspace{-0.05cm}
  \begin{block}{Objetivo de la Sesión}
    Derivar rigurosamente las fórmulas de Mínimos Cuadrados Ordinarios (MCO) mediante cálculo diferencial, demostrar la descomposición fundamental $\text{SCT}=\text{SCR}+\text{SCE}$, y establecer las propiedades del coeficiente de determinación $R^2$.
  \end{block}
\end{frame}

% Slide 3: Motivación
\begin{frame}{El Objetivo: Minimizar el Error de Ajuste}
  \footnotesize
  \begin{motivacion}
    Proponemos el modelo $Y_e=\hat\alpha+\hat\beta X$ para estimar $Y$ a partir de $X$. El propósito de la regresión lineal es encontrar los valores $\hat\alpha,\hat\beta$ que \emph{minimicen} la diferencia entre $Y$ y $Y_e$, representada por el error $\epsilon=Y-Y_e$.
  \end{motivacion}

  \vspace{0.15cm}
  \pause
  El método de \textbf{mínimos cuadrados} minimiza la suma de los errores al cuadrado, no los errores absolutos: esto hace al problema diferenciable y penaliza más los errores grandes.
\end{frame}

% Slide 4: Función objetivo
\begin{frame}{La Función Objetivo de Mínimos Cuadrados}
  \footnotesize
  \begin{block}{Función Objetivo}
    $$S(\alpha,\beta) = \sum_{i=1}^n\left(Y_i-(\alpha+\beta X_i)\right)^2 \quad \longrightarrow \quad \min_{\alpha,\beta} S$$
  \end{block}
  \pause

  \vspace{0.1cm}
  \begin{alertblock}{Estrategia de Minimización}
    Igualamos a cero las derivadas parciales $\partial S/\partial\alpha=0$ y $\partial S/\partial\beta=0$: el sistema resultante se conoce como las \textbf{ecuaciones normales}.
  \end{alertblock}
\end{frame}

% Slide 5: Derivación
\begin{frame}{Derivación de las Fórmulas de Mínimos Cuadrados}
  \footnotesize
  \begin{block}{Derivada respecto a $\alpha$}
    $$\frac{\partial S}{\partial\alpha}=-2\sum_i(Y_i-\alpha-\beta X_i)=0 \implies \hat\alpha=\bar Y-\hat\beta\bar X$$
  \end{block}
  \pause

  \vspace{0.1cm}
  \begin{alertblock}{Derivada respecto a $\beta$ (sustituyendo $\hat\alpha$)}
    $$\hat\beta=\frac{\sum_i(X_i-\bar X)(Y_i-\bar Y)}{\sum_i(X_i-\bar X)^2}=\frac{S_{XY}}{S_{XX}}$$
  \end{alertblock}
\end{frame}

% Slide 6: Teorema de partición
\begin{frame}{Teorema de Partición: SCT = SCR + SCE}
  \footnotesize
  \begin{block}{Identidad de Partición}
    $$\underbrace{\sum(Y_i-\bar Y)^2}_{\text{SCT}} = \underbrace{\sum(\hat Y_i-\bar Y)^2}_{\text{SCR}} + \underbrace{\sum(Y_i-\hat Y_i)^2}_{\text{SCE}}$$
  \end{block}
  \pause

  \vspace{0.1cm}
  \begin{alertblock}{Clave de la Demostración}
    El término cruzado $\sum(\hat Y_i-\bar Y)(Y_i-\hat Y_i)$ se anula exactamente porque las ecuaciones normales garantizan $\sum e_i=0$ y $\sum X_ie_i=0$ (y por tanto $\sum\hat Y_ie_i=0$).
  \end{alertblock}
\end{frame}

% Slide 7: R²
\begin{frame}{El Coeficiente de Determinación $R^2$}
  \footnotesize
  \begin{block}{Definición y Propiedades}
    $$R^2=\frac{\text{SCR}}{\text{SCT}}=1-\frac{\text{SCE}}{\text{SCT}} \in[0,1], \qquad R^2=1 \iff \text{SCE}=0$$
  \end{block}
  \pause

  \vspace{0.1cm}
  \begin{alertblock}{$R^2$ Ajustado}
    $$R^2_{\text{adj}}=1-\frac{n-1}{n-p-1}(1-R^2)$$
    $R^2$ nunca disminuye al agregar predictores; $R^2_{\text{adj}}$ sí puede disminuir si el nuevo predictor no aporta suficiente poder explicativo.
  \end{alertblock}
\end{frame}

% Slide 8: Lab Python (Bloque 1)
\begin{frame}[fragile]{Laboratorio en Python: Verificación de las Ecuaciones Normales}
  \scriptsize
  Confirmación numérica de que MCO satisface $\sum e_i=0$ y $\sum X_ie_i=0$ (\texttt{09.03\_mathematics\_of\_regression.py}):
  {\lstset{basicstyle=\fontsize{3.6pt}{4.3pt}\selectfont\ttfamily,aboveskip=0.05em,belowskip=0.05em}
  \lstinputlisting[firstline=17, lastline=27]{../../code/09_regresiones/09.03_mathematics_of_regression.py}}
\end{frame}

% Slide 9: Lab Python (Bloque 2)
\begin{frame}[fragile]{Laboratorio en Python: Descomposición SCT=SCR+SCE}
  \scriptsize
  Verificación de la identidad y del término cruzado nulo que la demuestra:
  {\lstset{basicstyle=\fontsize{3.4pt}{4.1pt}\selectfont\ttfamily,aboveskip=0.05em,belowskip=0.05em}
  \lstinputlisting[firstline=32, lastline=44]{../../code/09_regresiones/09.03_mathematics_of_regression.py}}
\end{frame}

% Slide 10: Lab Python (Bloque 3)
\begin{frame}[fragile]{Laboratorio en Python: Propiedades de $R^2$}
  \scriptsize
  Identidad $R^2=r^2$ y comportamiento del $R^2$ ajustado ante un predictor inútil:
  {\lstset{basicstyle=\fontsize{3.2pt}{3.9pt}\selectfont\ttfamily,aboveskip=0.05em,belowskip=0.05em}
  \lstinputlisting[firstline=48, lastline=69]{../../code/09_regresiones/09.03_mathematics_of_regression.py}}
\end{frame}

% Slide 11: Salida Terminal
\begin{frame}[fragile]{Laboratorio en Python: Salida en Terminal}
  \scriptsize
  Salida estándar generada al ejecutar el script del laboratorio en Python:
  \vspace{0.02cm}
  \begin{block}{Salida Estándar en Consola}
    \fontsize{3.6pt}{4.3pt}\selectfont
    \begin{verbatim}
=== Block 1: Normal Equations ===
sum(residuals)=0.0000000000, sum(x*residuals)=0.0000000000

=== Block 2: SST=SSR+SSE ===
SST=14309.41, SSR=12304.00, SSE=2005.41; cross term=0.0000000000

=== Block 3: R^2 Properties ===
R^2=0.8599 = r^2=0.8599
Adding useless predictor: R^2 0.8599->0.8602 (up); adj-R^2 0.8591->0.8587 (down)
    \end{verbatim}
  \end{block}
\end{frame}

% Slide 15: Cierre
\begin{frame}{Cierre de Sección: Matemáticas de la Regresión}
  \begin{reflexion}
    Hemos derivado rigurosamente las fórmulas MCO desde primeros principios (cálculo diferencial), demostrado que la descomposición SCT=SCR+SCE es una consecuencia exacta de las ecuaciones normales, y establecido las propiedades fundamentales del coeficiente $R^2$.
  \end{reflexion}

  \vspace{0.2cm}
  \pause
  \begin{block}{Lecciones Clave}
    \begin{itemize}\itemsep=0.08cm
      \item \textbf{MCO:} $\hat\beta=S_{XY}/S_{XX}$, $\hat\alpha=\bar Y-\hat\beta\bar X$, insesgados por construcción.
      \item \textbf{Ortogonalidad:} $\sum e_i=0$ y $\sum X_ie_i=0$ son la clave algebraica de toda la teoría posterior.
      \item \textbf{$R^2$ vs. $R^2_{\text{adj}}$:} el primero nunca disminuye; el segundo penaliza predictores poco útiles.
    \end{itemize}
  \end{block}
\end{frame}

% Slide 16: Perspectiva Modular
\begin{frame}{Perspectiva Modular --- El Próximo Reto}
  \scriptsize
  \begin{retoclase}
    Con las fórmulas derivadas rigurosamente, la Sección 09.04 las aplica a un conjunto de datos simulados generado explícitamente con parámetros conocidos ($\alpha,\beta,\sigma$ verdaderos), permitiendo comparar directamente los valores estimados por MCO contra los valores poblacionales reales que generaron los datos.
  \end{retoclase}

  \vspace{0.08cm}
  \pause
  \begin{alertblock}{Preparación Recomendada}
    \begin{itemize}\itemsep=0.06cm
      \item Repasa la diferencia entre parámetros poblacionales ($\alpha,\beta$) y sus estimadores muestrales ($\hat\alpha,\hat\beta$).
      \item Reflexiona sobre qué tan cerca esperarías que $\hat\beta$ esté de $\beta$ al aumentar el tamaño de muestra $n$.
    \end{itemize}
  \end{alertblock}
\end{frame}

\end{document}
